\documentclass{article}
\usepackage[utf8]{inputenc}
\usepackage[T1]{fontenc}
\usepackage{geometry}
\usepackage{graphicx}
\usepackage{float}
\usepackage{hyperref}
\usepackage[english,indonesian]{babel}

\title{Optimasi Pemrosesan Template LaTeX Menggunakan Pendekatan Berbasis AI}
\author{Academic Support Agent}
\date{\today}

\begin{document}
\maketitle

\begin{abstract}

\end{abstract}

\chapter{PENDAHULUAN}


\chapter{TINJAUAN PUSTAKA}


% Original content (commented out):
% %------------------------------------------------------------------------------
% SIMPLE TEMPLATE FOR P1 SCIENTIFIC JUSTIFICATION
% Use this template to generate the "scientific justification"
% for an ESO proposal for observing time, to be uploaded at
% https://www.eso.org/p1
%
% This LaTeX2e template is to be compiled with pdflatex.
% Import the resulting pdf file in the P1 interface.
% This file can be uploaded in a collaborative edition system like 
% OverLeaf.
%
% Do not tamper with the settings; changing the font size,
% margins, etc... may result in the proposal not being 
% considered.
% 
\documentclass{article}
\renewcommand{\thesection}{}
\renewcommand{\thesubsection}{}
\renewcommand{\familydefault}{\sfdefault}
\paperwidth 21cm
\paperheight 29.7cm
\textwidth 16.7cm
\textheight25.7cm
\oddsidemargin0cm \evensidemargin0cm
\topmargin-2cm
\parindent0cm
\parskip0.1cm
\usepackage{graphicx}
\pagestyle{empty}
\begin{document}
%------------------------------------------------------------------------------
\section{DESCRIPTION OF THE PROPOSED PROGRAMME}
\subsection{A- Scientific Rationale}

This template is for Normal, DDT and Large Proposals. Replace this text by your scientific justification. 

For Normal and DDT Proposals, the full text must be 2 pages or less, including figures. Page 3 and following will be cut. For Large Proposals, the text can be up to 3 pages, with 2 additional pages only for figures. Page 6 and following will be cut. When ready, process with pdflatex, and import the PDF into the P1 interface. 

The \LaTeX\ settings in the header should not be modified, and the font size or style should not be modified for the main text.\\


{\bf Scientific rationale: in this section, discuss the scientific background of  the project, pertinent references; previous work plus justification for present proposal}. This text can be removed.

Lorem ipsum dolor sit amet, consectetur adipiscing elit. Proin sit amet rhoncus nunc. Integer accumsan aliquam ipsum. Sed et purus pharetra, luctus enim ut, tempor ex. Pellentesque pharetra euismod cursus. Donec a est nisl. Nulla facilisi. Cras quis leo et lorem fermentum bibendum. Nunc volutpat, libero tristique molestie ornare, elit lorem venenatis magna, at tincidunt nulla est sed ante. Curabitur at dolor et nulla sagittis facilisis eget et nibh.

Praesent interdum augue ut ornare placerat. Maecenas cursus, augue eu aliquet gravida, sapien magna euismod leo, eu mollis sapien est quis justo. Vestibulum in lacus sed arcu vulputate porttitor. Sed commodo sodales orci sit amet consequat. In elementum in metus vel scelerisque. Donec libero magna, malesuada vitae felis ac, egestas semper diam. Donec eu velit dolor. Duis suscipit lobortis massa non efficitur. Vivamus condimentum turpis tempor massa consequat, in ullamcorper tortor volutpat. Class aptent taciti sociosqu ad litora torquent per conubia nostra, per inceptos himenaeos. Mauris finibus mauris lorem, id suscipit lectus imperdiet in. Etiam lectus lectus, dapibus at augue sit amet, cursus vestibulum odio. Cras commodo fermentum facilisis. Sed et convallis leo.

Mauris rhoncus porttitor sagittis. Sed pulvinar sollicitudin mollis. Vestibulum id aliquam mauris, eu malesuada libero. Quisque quis dui eros. Sed vitae eros nisl. Cras luctus, leo non volutpat dapibus, enim urna pulvinar ex, sed molestie justo erat in massa. Orci varius natoque penatibus et magnis dis parturient montes, nascetur ridiculus mus. Nunc et viverra nisi, ac rhoncus urna.

Phasellus sodales iaculis felis, ut elementum nibh tempus sit amet. Aenean ultricies nisl a nunc rutrum pellentesque. Nullam bibendum, ligula a auctor sagittis, lectus ligula vulputate metus, in accumsan nulla enim nec quam. Sed porttitor efficitur facilisis. Aliquam sed elementum metus, id iaculis lorem. In mollis elit ut lacus ultricies efficitur. Donec in tortor libero. Nunc id turpis at dui efficitur ornare vel vel lectus.

Vestibulum convallis odio et felis fermentum, id porta augue pellentesque. Vivamus velit mauris, pharetra eu aliquam vel, tincidunt luctus felis. Aenean accumsan vulputate risus, quis mollis est pulvinar consequat. Curabitur suscipit urna eget nisi scelerisque venenatis. Sed sagittis elementum massa quis mollis. Etiam ut pretium tortor. Duis vel ante ante.

Fusce volutpat viverra ex, vel rhoncus mauris imperdiet in. Vestibulum pharetra malesuada neque, ac elementum tortor. Aliquam augue neque, ornare eget bibendum eu, pharetra eget nulla. Donec bibendum suscipit nulla eu rutrum. Nunc id varius nunc, lacinia vehicula nulla. Proin sodales faucibus massa, eu bibendum tellus gravida in. Cras risus quam, pharetra sit amet lacinia non, tempor eu libero. Quisque ac accumsan ipsum. Integer congue finibus risus, nec rhoncus est aliquet eget. Etiam ac libero id orci sagittis pretium scelerisque nec nulla. Cras urna diam, efficitur ut neque sed, vehicula tincidunt justo. Vestibulum quis tortor et justo molestie accumsan. Pellentesque lectus elit, mattis nec dolor fringilla, lacinia aliquet urna. Etiam dolor tortor, pulvinar eu dolor a, lacinia blandit nisi.

\subsection{B- Immediate Objective}
{\bf Immediate objective of the proposal: state what is actually going to be observed and what shall be extracted from the observations, so that the feasibility becomes clear.} This text can be removed.

Lorem ipsum dolor sit amet, consectetur adipiscing elit. Proin sit amet rhoncus nunc. Integer accumsan aliquam ipsum. Sed et purus pharetra, luctus enim ut, tempor ex. Pellentesque pharetra euismod cursus. Donec a est nisl. Nulla facilisi. Cras quis leo et lorem fermentum bibendum. Nunc volutpat, libero tristique molestie ornare, elit lorem venenatis magna, at tincidunt nulla est sed ante. Curabitur at dolor et nulla sagittis facilisis eget et nibh.

Praesent interdum augue ut ornare placerat. Maecenas cursus, augue eu aliquet gravida, sapien magna euismod leo, eu mollis sapien est quis justo. Vestibulum in lacus sed arcu vulputate porttitor. Sed commodo sodales orci sit amet consequat. In elementum in metus vel scelerisque. Donec libero magna, malesuada vitae felis ac, egestas semper diam. Class aptent taciti sociosqu ad litora torquent per conubia nostra, per inceptos himenaeos. Mauris finibus mauris lorem, id suscipit lectus imperdiet in. Etiam lectus lectus, cursus vestibulum odio. Sed et convallis leo.

Proin venenatis et eros at sodales. Mauris vitae lacinia risus. Donec velit purus, placerat eu aliquam a, maximus ut nisi. Quisque euismod, ipsum id porttitor tincidunt, dui mi maximus purus, sit amet consectetur libero tortor ac nibh. Aenean mattis efficitur quam, sit amet porta dui sollicitudin in. Etiam porttitor ante non quam pulvinar, sit amet tempor odio mollis: $i$ Mauris rhoncus porttitor sagittis. $ii$ Sed pulvinar sollicitudin mollis. Vestibulum id aliquam mauris, eu malesuada libero. Quisque quis dui eros. $iii$ Sed vitae eros nisl. Cras luctus, leo non volutpat dapibus, enim urna pulvinar ex, sed molestie justo erat in massa. Orci varius natoque penatibus et magnis dis parturient montes, nascetur ridiculus mus. Nunc et viverra nisi, ac rhoncus urna.

Phasellus sodales iaculis felis, ut elementum nibh tempus sit amet. Aenean ultricies nisl a nunc rutrum pellentesque. Nullam bibendum, ligula a auctor sagittis, lectus ligula vulputate metus, in accumsan nulla enim nec quam. Sed porttitor efficitur facilisis. Aliquam sed elementum metus, id iaculis lorem. In mollis elit ut lacus ultricies efficitur. Donec in tortor libero. Nunc id turpis at dui efficitur ornare vel vel lectus.

Fusce velit metus, laoreet nec viverra ac, sollicitudin viverra augue. Vivamus dignissim rhoncus lorem ac sollicitudin. Praesent consequat ipsum ac nunc luctus sollicitudin a et metus. Nam consequat lacus quis diam venenatis scelerisque. Proin dapibus urna ullamcorper massa porttitor, id semper metus venenatis. Interdum et malesuada fames ac ante ipsum primis in faucibus. Mauris eget congue odio, quis placerat urna.

\subsection{Figures}
\includegraphics[width=0.48\textwidth]{Fig1.pdf}
\includegraphics[width=0.48\textwidth]{Fig2.pdf}

Fig.1 (left) This is the caption of the first figure. This plot is just an example illustrating the power of Python's Matplotlib package. Fig.2 (right). Feel free to use a Table to structure the plots and captions. Credit: matplotlib.org

\subsection{References}

The references can use a smaller font (eg 8pt,or footnotesize). What follows is an example, not a strict guideline. A suitable format can be generated by ADS using the following custom format: 
\begin{verbatim} %z0%3i (%Y) "%\T", %\q, %\V, %\p 
\end{verbatim}

{\footnotesize
{\bf 1.} Kuijken, K., et al. (2015) "Gravitational lensing analysis of the Kilo-Degree Survey," MNRAS,  454,  3500  -- 
{\bf 2.} de Jong, J. T. A., et al. (2015) "The first and second data releases of the Kilo-Degree Survey," A\&A,  582,  A62 -- 
{\bf 3.} Sobral, D., et al. (2015) "Evidence for PopIII-like Stellar Populations in the Most Luminous Lyman-$\alpha$ Emitters at the Epoch of Reionization: Spectroscopic Confirmation," ApJ,  808,  139 -- 
{\bf 4.} Massey, R., et al. (2015) "The behaviour of dark matter associated with four bright cluster galaxies in the 10 kpc core of Abell 3827," MNRAS,  449,  3393 -- 
{\bf 5.} Le Fevre, O., et al. (2015) "The VIMOS Ultra-Deep Survey: ~10 000 galaxies with spectroscopic redshifts to study galaxy assembly at early epochs $2 < z < 6$," A\&A,  576,  A79 -- 
{\bf 6.} Smette, A., et al. (2015) "Molecfit: A general tool for telluric absorption correction. I. Method and application to ESO instruments," A\&A,  576,  A77 -- 
{\bf 7.} Grazian, A., et al. (2015) "The galaxy stellar mass function at $3.5 < z < 7.5$ in the CANDELS/UDS, GOODS-South, and HUDF fields," A\&A,  575,  A96 -- 
{\bf 8.} Bacon, R., et al. (2015) "The MUSE 3D view of the Hubble Deep Field South," A\&A,  575,  A75 
}

\end{document}


\end{document}